\documentclass[11pt,a4paper]{article}

\usepackage{amsmath,amssymb,amsthm}
\usepackage{hyperref}
\usepackage[margin=1in]{geometry}
\usepackage{xcolor}
\usepackage{booktabs}
\usepackage{array}
\usepackage{float}
\usepackage{mathtools}

\newtheorem{theorem}{Theorem}[section]
\newtheorem{conjecture}[theorem]{Conjecture}
\newtheorem{corollary}[theorem]{Corollary}
\newtheorem{lemma}[theorem]{Lemma}
\newtheorem{proposition}[theorem]{Proposition}
\newtheorem{definition}[theorem]{Definition}
\theoremstyle{remark}
\newtheorem{remark}[theorem]{Remark}

\newcommand{\PT}{\mathbb{PT}}
\newcommand{\CP}{\mathbb{CP}}
\newcommand{\C}{\mathbb{C}}
\newcommand{\Z}{\mathbb{Z}}
\newcommand{\R}{\mathbb{R}}
\newcommand{\F}{\mathbb{F}}
\newcommand{\cO}{\mathcal{O}}
\newcommand{\cC}{\mathcal{C}}
\newcommand{\PP}{\mathcal{P}}
\newcommand{\wtSp}{\widetilde{Sp}}
\newcommand{\wtG}{\widetilde{G}}


\title{The Weil Representation and the $S_3$ Lifting Classification\\
       for Mermin Pentagrams\\
       {\large (Paper XV)}}

\author{Zhou-Li Chen \\ co-nlang Research}
\date{June 2026}

\begin{document}

\maketitle

\begin{abstract}
We apply the Weil (metaplectic) representation of $Sp(6,\mathbb{F}_2)$ on
$\mathbb{C}^8$ to the orbit structure of Mermin pentagrams established in
Papers~XI--XIV.
Our main results are:
(1)~the \emph{$S_3$ Lifting Classification} (Theorem~\ref{thm:s3lift}):
the three non-conjugate $S_3$ stabilizer subgroups of $G_2(2)$ lift
differently to the metaplectic group---$O_3$ lifts as a \emph{split} extension
($\widetilde{S}_3 \cong S_3$), $O_2$ as a \emph{non-split symmetric} extension
(both order-3 generators acquire metaplectic order~6), and $O_4$ as a
\emph{non-split asymmetric} extension (the two order-3 generators acquire
different metaplectic orders);
(2)~a \emph{new orbit structure}: $G_2(2)$ acts on the 135 Lagrangians of
$W(5,\mathbb{F}_2)$ in exactly two orbits of sizes 72 and 63, with
$\mathrm{Stab}_{G_2(2)}(V) \cong GL(3,\mathbb{F}_2)$;
and (3)~a \emph{negative result}: the naive conjecture
$\beta(C)/2 \equiv c(g_C, g_C^{-1}) \pmod 2$ fails (48.8\% match),
with a structural diagnosis showing that $\beta(C)$ is a \emph{context-level}
invariant (depending on the specific 4 rays of $C$) rather than a
Lagrangian-level invariant---the correct metaplectic identification of
$\beta$ remains an open problem.
\end{abstract}


\section{Introduction}

Papers~XI through~XIV established the parity theorem for Mermin pentagrams
and its orbit-theoretic consequences \cite{PaperXI,PaperXII,PaperXIII,PaperXIV}.
The central mystery is the \emph{mod~4 lifting problem}: the integer symplectic
form $\omega_{\mathrm{int}}$ gives $\beta_{\mathrm{sum}} \equiv 2 \pmod 4$ for all
12{,}096 pentagrams, but existing mod-2 theory only forces $\beta_{\mathrm{sum}}
\equiv 0 \pmod 2$.
Paper~XIV identified three open problems pointing toward the metaplectic framework:
the Orbit~4 parity anomaly (60.6\% 1-minus vs.\ 65--68\%), the three non-conjugate
$S_3$ stabilizers of $G_2(2)$, and the 56-multiple structure of all
orbit~$\times$~$k$-profile counts.

This paper applies the Weil representation $\rho: G_2(2) \to GL(8,\mathbb{C})$
(restriction of the Weil representation of $Sp(6,\mathbb{F}_2)$) to these problems.
Our main results are:

\begin{enumerate}
  \item \textbf{$S_3$ Lifting Classification} (\S\ref{sec:s3lift}):
        Three distinct lifting types for the three orbit stabilizers;
        $O_3$ is the unique split orbit.

  \item \textbf{$G_2(2)$ Lagrangian Orbit Structure} (\S\ref{sec:lagorbits}):
        $G_2(2)$ has exactly two orbits on the 135 Lagrangians (sizes 72 and 63);
        the stabilizer of $V$ is $GL(3,\mathbb{F}_2)$.

  \item \textbf{$\beta$-Cocycle Negative Result} (\S\ref{sec:betaneg}):
        The Lagrangian-level cocycle identification fails;
        $\beta(C)$ is a context-level invariant.
\end{enumerate}

\noindent
All computations use the scripts \texttt{metaplectic\_lift.py} and
\texttt{weil\_cocycle.py} in \cite{PaperXVsuppl}, running in approximately
60 seconds on $G_2(2)$.


\section{\texorpdfstring{Weil Representation of $Sp(6,\mathbb{F}_2)$}{Weil Representation of Sp(6,F2)}}\label{sec:weil}

\subsection{\texorpdfstring{Construction on $\mathbb{C}^8$}{Construction on C8}}

The Weil (oscillator) representation of $Sp(2n,\mathbb{F}_q)$ acts on
$L^2(\mathbb{F}_q^n) \cong \mathbb{C}^{q^n}$.
For $n=3$, $q=2$, the representation space is $\mathbb{C}^8$, with
basis $\{|x\rangle : x \in \mathbb{F}_2^3\}$.

$Sp(6,\mathbb{F}_2)$ is generated by three types of elements
\cite{Gerardin, GurevichHadani}:

\begin{enumerate}
  \item \textbf{Linear} ($A \in GL(3,\mathbb{F}_2)$):
        $\rho\!\left(\left[\begin{smallmatrix}A & 0 \\ 0 & A^{-T}\end{smallmatrix}\right]\right)
        |x\rangle = |Ax\rangle$.

  \item \textbf{Symmetric} ($B = B^T \in M_{3\times 3}(\mathbb{F}_2)$):
        $\rho\!\left(\left[\begin{smallmatrix}I & B \\ 0 & I\end{smallmatrix}\right]\right)
        |x\rangle = (-1)^{x^T B x}|x\rangle$.

  \item \textbf{Fourier}:
        $\rho\!\left(\left[\begin{smallmatrix}0 & I \\ I & 0\end{smallmatrix}\right]\right)
        |x\rangle = 2^{-3/2} \sum_{y \in \mathbb{F}_2^3} (-1)^{x \cdot y} |y\rangle$.
\end{enumerate}

Any $g \in Sp(6,\mathbb{F}_2)$ can be decomposed into these generators via
LDU decomposition (when the $D$-block is invertible) or via a Fourier
pre-multiplication, and $\rho(g)$ is computed as the product of the
corresponding matrices.

\begin{remark}[The $q=2$ case and Gauss sums]
For $q$ odd, the Weil cocycle is given by
$c(g,h) = \gamma(\mu(L_0, gL_0, ghL_0))$
where $\gamma$ is the Gauss sum and $\mu$ is the Kashiwara triple index.
For $q=2$, the Gauss sum $\sum_{x \in \mathbb{F}_2} (-1)^{x^2} = 1 + (-1) = 0$
degenerates, so this formula does not apply.
G\'erardin~\cite{Gerardin} constructs the Weil representation for all finite
fields including $q=2$ using a different method; the cocycle is extracted directly
from the representation matrices via $\rho(g)\rho(h) = c(g,h)\rho(gh)$.
We use this matrix-based extraction throughout.
\end{remark}

\subsection{\texorpdfstring{$G_2(2)$ and the Metaplectic Cocycle}{G2(2) and the Metaplectic Cocycle}}

We restrict $\rho$ to $G_2(2) \leq Sp(6,\mathbb{F}_2)$ (using the standard
symplectic basis, obtained from the ATLAS generators via change of basis).
The restriction is unitary: $\rho(g)\rho(g)^\dagger = I_8$ for all
$g \in G_2(2)$ (verified by spot-check).

For $g, h \in G_2(2)$, the \emph{metaplectic cocycle} is
\[
  c(g, h) = \mathrm{Tr}\bigl(\rho(g)\rho(h)\rho(gh)^{-1}\bigr)/8 \;\in\; \{+1, -1\}.
\]
Since $\rho$ is a projective representation with 2-cocycle values in $\{+1,-1\}$,
the product $\rho(g)\rho(h)\rho(gh)^{-1}$ is always a scalar matrix
$c(g,h)\cdot I_8$, so the trace formula directly extracts $c(g,h)$.
This satisfies the 2-cocycle condition $c(g,h)c(gh,k) = c(g,hk)c(h,k)$
and represents the central $\mathbb{Z}/2$ ambiguity in the metaplectic lift
$\tilde{g} \in \wtSp(6,\mathbb{F}_2)$ of $g \in Sp(6,\mathbb{F}_2)$.


\section{\texorpdfstring{$G_2(2)$ Lagrangian Orbit Structure}{G2(2) Lagrangian Orbit Structure}}\label{sec:lagorbits}

\begin{theorem}[$G_2(2)$ orbits on Lagrangians]\label{thm:lagorbits}
$G_2(2)$ acts on the 135 Lagrangians of $W(5,\mathbb{F}_2)$ with exactly two orbits:
\[
  \{\text{135 Lagrangians}\} = \mathrm{Orb}_A \;\sqcup\; \mathrm{Orb}_B,
  \quad |\mathrm{Orb}_A| = 72, \quad |\mathrm{Orb}_B| = 63.
\]
The standard Lagrangian $V = \mathbb{F}_2^3 \times \{0\}^3$ lies in $\mathrm{Orb}_A$,
with stabilizer
$\mathrm{Stab}_{G_2(2)}(V) \cong GL(3,\mathbb{F}_2)$ of order~$168 = 12{,}096/72$.
The complementary orbit $\mathrm{Orb}_B$ has stabilizer of order
$192 = 12{,}096/63 = 2^6 \times 3$.
\end{theorem}

\begin{remark}[Significance of $|\mathrm{Orb}_B| = 63$]
The number $63 = 2^6 - 1$ equals the number of non-zero vectors in $\mathbb{F}_2^6$,
or equivalently the number of points of $PG(5,\mathbb{F}_2)$.
The 63 Lagrangians in $\mathrm{Orb}_B$ may be related to the canonical
embedding of the 63-point set in the projective geometry of the Cayley hexagon;
we leave this identification to future work.
\end{remark}

\begin{remark}[$k$-profile invariance and the $V$-orbit]
\label{rem:v-clarify}
The orbit data ($|\mathrm{Orb}_A| = 72$, $V \in \mathrm{Orb}_A$) implies
there exist $g \in G_2(2)$ with $g(V) \neq V$.
\textbf{Correction to Paper~XIV:}
Paper~XIV Remark~2.3 justified $k$-profile invariance by claiming that
$G_2(2)$ preserves $V$; Theorem~\ref{thm:lagorbits} above shows this is
incorrect---$G_2(2)$ moves $V$ through a 72-element orbit.
The computational verification of $k$-profile invariance in Paper~XIV
stands, but the stated justification must be replaced.

The correct picture is as follows.
When $g$ moves $V$ to $g(V) \neq V$, individual $k$-values satisfy
$|g(L) \cap V| = |L \cap g^{-1}V| \neq |L \cap V|$ in general.
However, $k$-profile invariance is a statement about the \emph{sorted tuple}
$\kappa(\mathbf{P}) = (k_1 \geq \cdots \geq k_5)$, not individual values.
Even when $g$ moves $V$, it acts on the pentagram's five Lagrangians
simultaneously; the sorted multiset of the five resulting $k$-values
(relative to the fixed $V$) is empirically unchanged.
Paper~XIV verified this for all 12{,}096 pentagrams
\cite{PaperXIVsuppl}; the algebraic reason is an open problem
(Problem~\ref{sec:open}, item~3).
\end{remark}


\section{\texorpdfstring{$S_3$ Lifting Classification}{S3 Lifting Classification}}\label{sec:s3lift}

\subsection{Setup}

Paper~XIV established that the four $G_2(2)$ orbits have stabilizers
$\mathrm{Stab}(O_1) \cong \mathbb{Z}_2$ and
$\mathrm{Stab}(O_i) \cong S_3$ for $i=2,3,4$,
with the three $S_3$ subgroups pairwise non-conjugate in $G_2(2)$
\cite{PaperXIV}.
We now classify how these $S_3$ subgroups lift to the metaplectic group via
$\rho$.

\subsection{Metaplectic Lift Invariants}

For $g \in S_3 \leq G_2(2)$, the Weil representation gives
$\rho(g) \in GL(8,\mathbb{C})$.
We compute three invariants:

\begin{itemize}
  \item \textbf{$s^2$-cocycle}: for each order-2 element $s$,
        the sign $c(s,s) \in \{+1,-1\}$, where $c(s,s) = +1$ iff
        $\rho(s)^2 = I$ (split) and $-1$ iff $\rho(s)^2 = -I$ (non-split,
        $\rho(s)$ has order~4 rather than~2).

  \item \textbf{$t^3$-cocycle}: for each order-3 element $t$,
        the sign $c(t,t) \cdot c(t^2,t) \in \{+1,-1\}$, where $+1$ means
        $\rho(t)$ has order~3 (split) and $-1$ means order~6 (non-split).

  \item \textbf{$(st)^2$-cocycle}: for each pair $(s,t)$ with $s$ order-2
        and $t$ order-3, the sign $c(st, st) \in \{+1,-1\}$.
        In $S_3$, $(st)^2 = e$; in the metaplectic lift, $(st)^2 = \pm I$.
\end{itemize}

\subsection{Main Theorem}

\begin{theorem}[$S_3$ Lifting Classification]\label{thm:s3lift}
The three non-conjugate $S_3$ stabilizer subgroups of $G_2(2)$ lift to three
distinct types in the metaplectic group.
Verification on one representative per orbit is sufficient: since the Weil
representation is $G_2(2)$-equivariant, any two elements $H_1, H_2 \leq G_2(2)$
related by $H_2 = gH_1 g^{-1}$ satisfy $\rho(H_2) = \rho(g)\rho(H_1)\rho(g)^{-1}$,
so conjugate stabilizers have the same lifting type.
\end{theorem}

\begin{table}[H]
\centering
\caption{Metaplectic lift data for the three order-6 orbit stabilizers.
$S_3 = \langle s, t \mid s^2 = t^3 = (st)^2 = e\rangle$.
Computed via the Weil representation on $\mathbb{C}^8$
\cite{PaperXVsuppl}.}
\label{tab:s3lift}
\begin{tabular}{lcccl}
\toprule
Orbit & $s^2$-cocycles & $t^3$-cocycles & $(st)^2 = +I$: & Lift type \\
\midrule
$O_3$ & $[+1, +1, +1]$ & $[+1, +1]$ & 6 of 6 & SPLIT \\
$O_2$ & $[+1, +1, -1]$ & $[-1, -1]$ & 4 of 6 & NON-SPLIT symmetric \\
$O_4$ & $[+1, +1, -1]$ & $[-1, +1]$ & 4 of 6 & NON-SPLIT asymmetric \\
\bottomrule
\end{tabular}
\end{table}

\noindent
The three lifting types are:
\begin{itemize}
  \item \textbf{SPLIT} ($O_3$): All three $s^2$-cocycles are $+1$ and both
        $t^3$-cocycles are $+1$.
        The central extension $1 \to \mathbb{Z}/2 \to \wtSp \to Sp \to 1$
        splits when restricted to the $O_3$ stabilizer: the lifted subgroup
        satisfies $\widetilde{S_3} \cong S_3$.

  \item \textbf{NON-SPLIT symmetric} ($O_2$): One $s^2$-cocycle is $-1$
        (one reflection has metaplectic order~4) and both $t^3$-cocycles
        are $-1$ (both order-3 elements have metaplectic order~6).
        The extension is non-trivial, but the two order-3 generators lift
        consistently.

  \item \textbf{NON-SPLIT asymmetric} ($O_4$): Same $s^2$-cocycle pattern
        as $O_2$, but the two $t^3$-cocycles differ: $[-1, +1]$.
        One order-3 generator has metaplectic order~6 while the other retains
        order~3.
\end{itemize}

\begin{remark}[Symmetry breaking in $O_4$]
In the abstract group $S_3$, the two order-3 elements $t$ and $t^{-1} = t^2$
are conjugate: $s \cdot t \cdot s = t^{-1}$ for any reflection $s$.
This inner automorphism makes $t$ and $t^{-1}$ indistinguishable in $S_3$.
The $O_4$ embedding in $G_2(2) \leq Sp(6,\mathbb{F}_2)$ breaks this symmetry:
$\rho(t)^3 = -I$ while $\rho(t^{-1})^3 = +I$ (or vice versa), so the two
conjugate elements land in different conjugacy classes of the metaplectic group.
This symmetry breaking is a property of the specific embedding of $S_3$ into
$G_2(2)$, invisible at the abstract group level.
\end{remark}

\begin{remark}[Three-level invariant hierarchy]
\label{rem:3level}
The three $G_2(2)$ orbits $O_2, O_3, O_4$ admit a three-level classification:
\end{remark}

\begin{table}[H]
\centering
\caption{Three-level classification of $G_2(2)$ orbits.
The metaplectic lifting level is the finest invariant, distinguishing
$O_2$ from $O_3$ within Class~I and providing a structural difference
between $O_2$ and $O_4$.}
\label{tab:3level}
\begin{tabular}{lccc}
\toprule
Level & $O_2$ & $O_3$ & $O_4$ \\
\midrule
Abstract stabilizer (Paper~XIV) & $S_3$ & $S_3$ & $S_3$ \\
Class I/II (Paper~X \cite{PaperX}) & Class I & Class I & Class II \\
Metaplectic lifting (this paper) & NS-symmetric & SPLIT & NS-asymmetric \\
\bottomrule
\end{tabular}
\end{table}

\begin{remark}[$O_3$ split lift and $k$-profile uniformity]
Orbit~$O_3$ is simultaneously the unique split orbit (Theorem~\ref{thm:s3lift})
and the unique orbit with equal proportions of the three dominant $k$-profiles:
$(1,1,1,0,0) = (1,0,0,0,0) = (3,1,1,0,0) = 25.0\%$
(Paper~XIV Table~A8, orbit $\times$ $k$-profile percentage cross-table
\cite{PaperXIVsuppl}).
Whether these two properties have a common group-theoretic origin is an open
question.
\end{remark}

\begin{remark}[Orbit~4 anomaly and the asymmetric lift]
\label{rem:orbit4asym}
The Orbit~4 parity anomaly (60.6\% 1-minus pentagrams vs.\ 65--68\% for $O_1$--$O_3$,
Paper~XIII \cite{PaperXIII}) was shown in Paper~XIV to be \emph{intra-profile}:
the orbit $\times$ $k$-profile cross-table shows that $O_4$ has a
\emph{smaller} fraction of the low-1-minus $(1,0,0,0,0)$ profile than other orbits
(16.7\% vs.\ 22--25\%), so $k$-profile composition alone would actually predict a
\emph{higher} 1-minus rate for $O_4$---the opposite of what is observed
\cite{PaperXIV,PaperXIVsuppl}.
The NON-SPLIT asymmetric lifting of $O_4$ provides the first structural
difference between $O_4$ and the other orbits at the group-theoretic level.

The mechanism hypothesis is:
the two order-3 generators of the $O_4$ stabilizer have different metaplectic
orders (6 and 3); their action on the parity function is correspondingly
asymmetric, shifting the balance between 1-minus and 3-minus pentagrams.
Quantifying this mechanism requires decomposing the parity function under the
$S_3$ representation on the parity fiber, which we leave to future work.
\end{remark}

\begin{remark}[$O_2$ and $O_4$ share the same $s^2$-cocycle pattern]
\label{rem:s2same}
From Table~\ref{tab:s3lift}, $O_2$ and $O_4$ have identical $s^2$-cocycle
vectors $[+1,+1,-1]$: in both orbits, two of the three reflections lift to
order-2 elements and one lifts to order-4.
Within the three order-6 orbits, the two types of cocycles play complementary roles:
$s^2$-cocycles distinguish $O_3$ from $\{O_2, O_4\}$ (all three are $+1$ for
$O_3$, while both $O_2$ and $O_4$ have one $-1$);
$t^3$-cocycles distinguish $O_2$ from $O_4$ ($[-1,-1]$ vs.\ $[-1,+1]$).
The full three-way classification requires both invariants.

For the Orbit~4 parity anomaly specifically---which concerns $O_4$ vs.\ the
other orbits---it is the $t^3$-cocycle asymmetry that is structurally responsible:
$O_4$'s two order-3 generators lift to different metaplectic orders, whereas
$O_2$'s lift symmetrically and $O_3$'s split entirely.
Consequently, the \emph{3-fold rotation} structure---not the
\emph{reflection} structure---carries the anomaly signal at the metaplectic level.
This may guide the context-level identification of $\beta$ in Paper~XVI:
the relevant symplectic transformation should capture 3-fold rotation
information, not 2-fold reflection.
\end{remark}


\section{\texorpdfstring{$\beta$-Cocycle Negative Result}{beta-Cocycle Negative Result}}\label{sec:betaneg}

\subsection{The Conjecture and Its Failure}

Based on the analogy with the real Maslov cocycle \cite{LionVergne1980},
we tested the identification:
\[
  \beta(C)/2 \;\equiv\; c(g_C,\, g_C^{-1}) \pmod 2,
\]
where $g_C \in Sp(6,\mathbb{F}_2)$ is any symplectic transformation with
$g_C(V) = L_C$ (mapping the standard Lagrangian $V$ to the Lagrangian $L_C$
containing context $C$).

\begin{theorem}[Lagrangian-level conjecture fails]\label{thm:betafails}
The identification $\beta(C)/2 \equiv c(g_C, g_C^{-1}) \pmod 2$ is false.
The match rate over all contexts in the 72 $G_2(2)$-accessible Lagrangians
is 48.8\%, consistent with random coincidence.
\end{theorem}

\subsection{Structural Diagnosis}

The failure is structural, not a computation artifact.
The key observation is that $\beta(C)$ depends on the specific \emph{context}
$C$ (which 4 rays from $L_C$), while $c(g_C, g_C^{-1})$ depends only on
the \emph{Lagrangian} $L_C$.

\begin{table}[H]
\centering
\caption{Fraction of Lagrangians with varying $h = \beta(C)/2 \bmod 2$
across their multiple contexts, by $k$-type.
``Varying'' means two contexts in the same Lagrangian have different $h$-values.
Verified over all 135 Lagrangians \cite{PaperXVsuppl}.}
\label{tab:lagvarying}
\begin{tabular}{lrrrr}
\toprule
$k$-type & \#Uniform & \#Varying & Total & \% varying \\
\midrule
$k = 0$ & 20 & 44 & 64 & 69\% \\
$k = 1$ & 24 & 32 & 56 & 57\% \\
$k = 3$ &  9 &  5 & 14 & 36\% \\
$k = 7$ &  1 &  0 &  1 & 0\% \\
\midrule
Total   & 54 & 81 & 135 & 60\% \\
\bottomrule
\end{tabular}
\end{table}

\noindent
Among the 135 Lagrangians, 81 (60\%) have at least two contexts with different
$h$-values.
Since $c(g_C, g_C^{-1})$ is constant across all contexts of the same Lagrangian,
it cannot match $\beta(C)/2$ for these 81 Lagrangians.

The $k=7$ case ($V$ itself) is the only exception: $V$ has $h = 0$ uniformly
for all its contexts (since all vectors in $V$ have $z$-component zero, so
$\omega_{\mathrm{int}} = 0$ and $\beta = 0$ for any context $C \subset V$),
consistent with $c(g_V, g_V^{-1}) = c(e, e) = 1$.

\begin{remark}[Context-level identification needed]
The failure of the Lagrangian-level conjecture suggests that the correct
metaplectic identification of $\beta(C)$ requires a symplectic transformation
$g_C$ that \emph{encodes the specific 4 rays of $C$}, not just the Lagrangian
$L_C$.
One natural candidate: the symplectic transformation associated to the
\emph{Clifford algebra element} or \emph{transvection} defined by the context.
This would make $g_C$ depend on $C$ intrinsically, resolving the
Lagrangian-level under-specification.
The correct context-level identification is the primary open problem for
Paper~XVI.
\end{remark}


\section{Open Problems}\label{sec:open}

\begin{enumerate}
  \item \textbf{Context-level metaplectic identification}:
        Find a symplectic transformation $\tilde{g}_C \in Sp(6,\mathbb{F}_2)$
        depending on the 4-ray context $C$ such that
        $\beta(C)/2 \equiv c(\tilde{g}_C, \tilde{g}_C^{-1}) \pmod 2$.
        Candidates: transvections, Clifford algebra elements, or the
        symplectic transformation defined by the basis of $C$.

  \item \textbf{Quantification of the $O_4$ anomaly via asymmetric lift}:
        Compute the action of the $O_4$ stabilizer $\widetilde{S}_3^{(O_4)}$
        on $\mathbb{C}^8$ and decompose by character.
        The asymmetric $t^3$-cocycle $[-1,+1]$ should produce an asymmetric
        measure distribution between the 1-minus and 3-minus parity sectors.
        Concretely: compute the character table of $\widetilde{S}_3^{(O_4)}$
        restricted to the parity fiber, and show that the resulting sector
        weights reproduce the observed 60.6\% 1-minus rate.
        Success would upgrade Remark~\ref{rem:orbit4asym} to a theorem.

  \item \textbf{Algebraic explanation of sorted $k$-profile invariance}:
        Theorem~\ref{thm:lagorbits} shows that $G_2(2)$ does \emph{not} fix
        $V$ setwise ($V \in \mathrm{Orb}_A$ with $|\mathrm{Orb}_A| = 72 > 1$),
        so the individual values $|g(L_i) \cap V|$ can change under
        $g \in G_2(2)$.
        Yet Paper~XIV verified computationally that the \emph{sorted} tuple
        $\kappa(\mathbf{P})$ is $G_2(2)$-invariant for all 12{,}096 pentagrams.
        The $\mathrm{Orb}_A$/$\mathrm{Orb}_B$ dichotomy alone cannot explain
        this: $\mathrm{Orb}_A$ contains Lagrangians of all four $k$-types
        ($k = 0, 1, 3, 7$), so orbit membership does not determine $k$-value.
        The correct explanation must account for how $G_2(2)$ simultaneously
        permutes a pentagram's five Lagrangians, preserving the multiset of
        $k$-values even as it moves $V$.
        A natural conjecture: the action of $G_2(2)$ on the five Lagrangians
        of any pentagram is a permutation that preserves the multiset
        $\{|L_i \cap V|\}$ because of the Cayley hexagon symmetry structure
        underlying $G_2(2)$.

  \item \textbf{Pentagram loop cocycle}:
        Once the context-level $\tilde{g}_C$ is identified, verify
        $\prod_{i=1}^5 c(\tilde{g}_{C_i}, \tilde{g}_{C_i}^{-1}) = -1$
        for all 12{,}096 pentagrams, proving
        $\beta_{\mathrm{sum}} \equiv 2 \pmod 4$ as a metaplectic 2-cocycle
        constraint.

  \item \textbf{$O_3$ split and $k$-profile uniformity}:
        Find a common group-theoretic explanation for the two coincidences
        of $O_3$: unique split lift and unique 25\%/25\%/25\% $k$-profile
        distribution.

  \item \textbf{$n$-qubit generalization}:
        Does the three-level classification (abstract stabilizer, orbit class,
        metaplectic lifting) extend to the $2n$-qubit Mermin pentagrams
        in $W(2n-1, \mathbb{F}_2)$?
\end{enumerate}


\appendix
\setcounter{table}{4}
\renewcommand{\thetable}{A\arabic{table}}
\renewcommand{\theHtable}{A\arabic{table}}
\section{Rigor Self-Assessment}

\begin{table}[H]
\centering
\caption{Rigor classification and verification scope for all main results.}
\label{tab:rigor}
\begin{tabular}{p{6.2cm}lp{4.5cm}}
\toprule
Result & Type & Scope \\
\midrule
Thm~\ref{thm:lagorbits} ($G_2(2)$ orbits on Lagrangians: $72+63$)
  & Computational & All 135 \cite{PaperXVsuppl} \\
Thm~\ref{thm:s3lift} ($S_3$ lifting classification)
  & Computational & One rep.\ per orbit \cite{PaperXVsuppl} \\
Rem~\ref{rem:orbit4asym} ($O_4$ anomaly via asymmetric lift)
  & Observational & Hypothesis; not quantified \\
Thm~\ref{thm:betafails} ($\beta$-cocycle conjecture refuted)
  & Algebraic refutation & 81/135 Lagrangians have varying $h$; Lagrangian-level is ill-defined \cite{PaperXVsuppl} \\
Table~\ref{tab:lagvarying} ($h$-variation per $k$-type)
  & Computational & All 135 Lagrangians \cite{PaperXVsuppl} \\
\bottomrule
\end{tabular}
\end{table}


\begin{thebibliography}{99}

\bibitem[PaperX]{PaperX}
Z.-L. Chen,
``The Equiangular Characterization of Mermin Pentagrams:
Symplectic Embedding, 10-Ray Cap, and the Failure of the
Collinearity--Obstruction Dictionary,''
\textit{Zenodo, DOI: 10.5281/zenodo.20476659} (2026).

\bibitem[PaperXI]{PaperXI}
Z.-L. Chen,
``Quadratic Refinement and the Parity of Mermin Pentagrams:
The $\beta$-Formula, Geometric Decomposition, and the
Kochen--Specker Obstruction from Equiangular Geometry,''
\textit{Zenodo, DOI: 10.5281/zenodo.20482595} (2026).

\bibitem[PaperXII]{PaperXII}
Z.-L. Chen,
``The T-Vector Theorem: Symplectic Characteristic Elements
and the Kochen--Specker Obstruction,''
\textit{Zenodo, DOI: 10.5281/zenodo.20490118} (2026).

\bibitem[PaperXIII]{PaperXIII}
Z.-L. Chen,
``The Maslov Index and the Kochen--Specker Obstruction:
Negative Results, the $k$-Profile Theorem, and Orbit~4 Anomaly,''
\textit{Zenodo, DOI: 10.5281/zenodo.20496513} (2026).

\bibitem[PaperXIV]{PaperXIV}
Z.-L. Chen,
``Stabilizer Algebra and the $k$-Profile Theorem for Mermin Pentagrams,''
\textit{Zenodo, DOI: 10.5281/zenodo.20501961} (2026).

\bibitem[PaperXIVsuppl]{PaperXIVsuppl}
Z.-L. Chen,
``Supplementary computational scripts for Paper~XIV,''
\textit{Zenodo, DOI: 10.5281/zenodo.20501961} (2026).

\bibitem[PaperXVsuppl]{PaperXVsuppl}
Z.-L. Chen,
``Supplementary computational scripts for Paper~XV,''
\textit{Zenodo, DOI: 10.5281/zenodo.20509690} (2026).

\bibitem[LionVergne1980]{LionVergne1980}
G. Lion and M. Vergne,
``The Weil Representation, Maslov Index and Theta Series,''
\textit{Birkh\"auser, Boston} (1980).

\bibitem[Gerardin]{Gerardin}
P. G\'erardin,
``Weil representations associated to finite fields,''
\textit{J. Algebra} \textbf{46} (1977), 54--101.

\bibitem[GurevichHadani]{GurevichHadani}
S. Gurevich and R. Hadani,
``Quantization of symplectic vector spaces over finite fields,''
\textit{J. Symplectic Geom.} \textbf{9} (2011), 1--12.

\end{thebibliography}

\end{document}
